On réalise une pile en utilisant le matériel et les produits suivants~:
\begin{itemize}
    \item un bêcher contenant le volume $V_{1}=\qty{20}{\mL}$ d'une solution
          aqueuse de nitrate d'argent $\ce{Ag^+(aq) + NO^-_3_{(aq)}}$ de
          concentration molaire $C_{1}=\qty{1.0e-1}{\mol\per\L}$
    \item un bêcher contenant le volume $V_{2}=\qty{20}{\mL}$ d'une solution
          aqueuse de nitrate de cuivre $\ce{Cu^2+_{(aq)} + 2NO^-_3_{(aq)}}$ de
          concentration molaire $C_{2}=\qty{5.0e-2}{\mol\per\L}$
    \item un fil de cuivre~;
    \item un fil d'argent~;
    \item un pont salin contenant une solution aqueuse saturée de nitrate de
          potassium $\ce{K^+(aq) + NO^-_3(aq)}$.
\end{itemize}

\begin{donnees}
    \begin{itemize}
        \item $1\mathcal{F}=\qty{96500}{\C\mol}$~;
        \item Constante d'équilibre associée à l'équation
              $\ce{2Ag^{+}_{(aq)} + Cu_{(s)} <=>[(1)][(2)] 2Ag_{(s)} +
              Cu^{2+}_{(aq)}}$ est~: $\mathrm{K}=\num{2.2e15}$
    \end{itemize}
\end{donnees}

On relie les électrodes de la pile à un conducteur ohmique en série avec un
ampèremètre, et on observe le passage d'un courant électrique dans le circuit
extérieur de la pile.
\begin{enumerate}
    \item Calculer la valeur du quotient de la réaction $Q_{r,i}$ dans l'état
          initial du système chimique. En déduire le sens spontané de
          l'évolution de ce système.
    \item On fait fonctionner la pile pendant une longue durée jusqu'à ce qu'il
          s'épuise. Déterminer la valeur de la quantité d'électricité qui
          traverse le conducteur ohmique depuis le début de fonctionnement de la
          pile jusqu'à son épuisement, sachant que le réactif limitant est l'ion
          $\ce{Ag^+}$.
\end{enumerate}



